% -----------------------------------------------------------
\section{Law}
% -----------------------------------------------------------
	\label{LAW}
	
	See also \hyperref[ANOMIA]{\grl{ἀνομία}}, \hyperref[PRIESTHOOD]{Priesthood}, \hyperref[COVENANT]{Covenant} (especially the \hyperref[COVENANT-TDOT]{TDOT} reference material),
	
% % ----------------------
% \subsection{See also}
% % ----------------------

% % ----------------------
% \subsection{1828 American Dictionary of the English Language} % *1828
% % ----------------------
	% \label{LAW-1828}

	% \begin{compactitem}
		% \item
	% \end{compactitem}

% % ----------------------
% \subsection{Oxford English Dictionary} % *OED
% % ----------------------
	% \label{LAW-OED}

	% \begin{compactitem}
		% \item[]
	% \end{compactitem}

% ----------------------
\subsection{Scriptural Usage} % *SCRIPTURES
% ----------------------
	\label{LAW-SCRIPTURES}

	% % -----
	% \subsubsection{Old Testament} % *OT
	% % -----
		% \label{LAW-OT}
	
		% % --
		% \paragraph{KJV}
		% % --
			% \label{LAW-OT-KJV}
	
			% \begin{tabular}{ll}
				% Total occurrences in the OT & 
			% \end{tabular}
			
			% \begin{compactdesc}
				% \item[]
			% \end{compactdesc}
			
		% % --
		% \paragraph{Hebrew Bible}
		% % --
			% \label{LAW-HEBREW}
		
			% %
			% \subparagraph{\texorpdfstring{\he{sample word}}{}}
			% %
				% \label{PAGA} % whatever is in step bible without periods
			
				% \begin{compactdesc}
					% \item[HALOT , PDF ] \textbf{}
						% \begin{compactitem}
							% \item[1]
						% \end{compactitem}
					% \item[BDB , PDF ] \textbf{}
						% \begin{compactitem}
							% \item[1]
						% \end{compactitem}
					% \item[DCH PDF ] \textbf{}
						% \begin{compactitem}
							% \item[1]
						% \end{compactitem}
				% \end{compactdesc}
				
				% \begin{compactdesc}
					% \item[Some OT uses] \textbf{}
						% \begin{compactdesc}
							% \item[]
						% \end{compactdesc}
				% \end{compactdesc}
			
		% % --
		% \paragraph{Septuagint}
		% % --
			% \label{LAW-LXX}
		
			% %
			% \subparagraph{\texorpdfstring{\gr{sample word}}{}}
			% %
		
	% -----
	\subsubsection{New Testament} % *NT
	% -----
		\label{LAW-NT}
	
		% --
		\paragraph{KJV}
		% --
			\label{LAW-NT-KJV}
			
	
			% \begin{tabular}{ll}
				% Total occurrences in the NT & 
			% \end{tabular}
			
			\begin{compactdesc}
				\item[{\hyperref[LUKE-16-16]{Luke 16:16--17}}]
				\item[{\hyperref[ROMANS-3]{Romans 3}}]
				\item[{\hyperref[ROMANS-4]{Romans 4}}] \phantom{.}
					\begin{compactitem}
						\item There's really material throughout Romans on this, clearly.
					\end{compactitem}
				\item[{\hyperref[ROMANS-7]{Romans 7--8}}]
				\item[{\hyperref[ROMANS-10-1]{Romans 10:1--10}}]
				\item[{\hyperref[ROMANS-13-8]{Romans 13:8--10}}]
				\item[1 Corinthians 15:56] The sting of death is sin; and the strength of sin is the law.
			\end{compactdesc}
			
		% --
		\paragraph{Greek New Testament} % *GREEK
		% --
			\label{LAW-GREEK}
		
			%
			\subparagraph{\texorpdfstring{\gr{ἀνόμως}}{}}
			%
				\label{ANOMOS-2}
			
				\begin{compactdesc}
					\item[BDAG 75a] \textbf{}
						\begin{compactitem}
							\item \textbf{without participation in an organized legal system, \textit{lawlessly}}
								\begin{compactitem}
									\item This entry glosses \hyperref[ROMANS-2-12]{Romans 2:12} as ``\textit{those who sin without law will also be lost without law}, i.e. those who sin without awareness of the Mosaic law (expressed from the vantage point of Israel's chosen status) will also perish without reference to it.''
									\item See also \hyperref[ROMANS-7-9]{Romans 7:9}.
								\end{compactitem}
						\end{compactitem}
					\item[LSJ , PDF ] \textbf{}
						\begin{compactitem}
							\item 
						\end{compactitem}
				\end{compactdesc}
				
				\begin{tabular}{ll}
					Total occurrences in the NT & 2
				\end{tabular}
				
				\begin{compactdesc}
					\item[Some NT uses] \textbf{}
						\begin{compactdesc}
							\item[{\hyperref[ROMANS-2-12]{Romans 2:12}}] \textbf{(2)}
						\end{compactdesc}
				\end{compactdesc}
		
	% -----
	\subsubsection{Book of Mormon} % *BOM
	% -----
		\label{LAW-BOM}
	
		% \begin{tabular}{ll}
			% Total occurrences in the BOM & 
		% \end{tabular}
		
		\begin{compactdesc}
			\item[{\hyperref[2NEPHI-9-17]{2 Nephi 9:17}}]
			\item[{\hyperref[ALMA-42]{Alma 42}}] \phantom{.}
				\begin{compactitem}
					\item Not every verse talks about laws, but the entire thing is important.
				\end{compactitem} 
			\item[{\hyperref[HELAMAN-6-24]{Helaman 6:24}}] \phantom{.}
				\begin{compactitem}
					\item Gives you an idea of what consequence-based laws for agents would be like in a system of law.
				\end{compactitem}
		\end{compactdesc}
		
	% -----
	\subsubsection{Doctrine and Covenants} % *DC
	% -----
		\label{LAW-DC}
	
		\begin{tabular}{ll}
			Total occurrences in the DC & 194
		\end{tabular}
		
		\begin{compactdesc}
			\item[{\hyperref[DC20-20]{DC 20:20}}] But by the transgression of these holy laws man became sensual and devilish, and became fallen man.
			\item[{\hyperref[DC29-34]{DC 29:34}}] Wherefore, verily I say unto you that all things unto me are spiritual, and not at any time have I given unto you a law which was temporal; neither any man, nor the children of men; neither Adam, your father, whom I created.
			\item[{\hyperref[DC38-21]{DC 38:21--22}}]
			\item[{\hyperref[DC38-32]{DC 38:32}}]
			\item[{\hyperref[DC41-3]{DC 41:3--5}}]
			\item[{\hyperref[DC41-10]{DC 41:10}}]
			\item[{\hyperref[DC42]{DC 42}}] \phantom{.}
				\begin{compactitem}
					\item Notice \hyperref[DC42-2]{verse 2} of this section---Christ delivers a specific law to them and tells them to obey. In this ``law'' stuff with the DC I keep thinking of \hyperref[2NEPHI-32-6]{2 Nephi 32:6}. See \hyperref[DC42-13]{DC 42:13} for example.)
					\item \hyperref[DC42-87]{Verse 87}
				\end{compactitem}
			\item[{\hyperref[DC43-8]{DC 43:8--9}}] \phantom{.}
				\begin{compactitem}
					\item Make sure to note the JSP versions, which are different.
				\end{compactitem}
			\item[{\hyperref[DC44-4]{DC 44:4--6}}]
			\item[{\hyperref[DC45-54]{DC 45:54}}]
			\item[{\hyperref[DC51]{DC 51}}] \phantom{.}
				\begin{compactitem}
					\item This section talks a lot about ``laws'' and the people organizing themselves according to the laws; it means specifically something like consecration.
				\end{compactitem}
			\item[{\hyperref[DC58-17]{DC 58:17--23}}]
			\item[{\hyperref[DC64-12]{DC 64:12--13}}]
			\item[{\hyperref[DC68-14]{DC 68:14--21}}] \phantom{.}
				\begin{compactitem}
					\item Notice the ``legal right'' stuff in talking about the priesthood; notice also that a lot of this was not in the original version of this revelation.
					\item Verses 16 and 20 are particularly important for thinking about laws.
				\end{compactitem}
			\item[{\hyperref[DC70-8]{DC 70:8}}]
			\item[{\hyperref[DC82-3]{DC 82:3--4}}]
			\item[{\hyperref[DC98-4]{DC 98:4--8}}]
			\item[{\hyperref[DC98-32]{DC 98:32--33}}]
			\item[{\hyperref[DC101-77]{DC 101:77--78}}]
			\item[{\hyperref[DC102-4]{DC 102:4}}]
			\item[{\hyperref[DC105-4]{DC 105:4--5}}]
			\item[{\hyperref[DC130-20]{DC 130:20--21}}]
			\item[{\hyperref[DC132-5]{DC 132:5}}]
			\item[{\hyperref[DC132]{DC 132}}] \phantom{.}
				\begin{compactitem}
					\item Lots of stuff in this section relevant to the priesthood and laws; too much to discuss here, at least for now.
				\end{compactitem}
		\end{compactdesc}
		
	% % -----
	% \subsubsection{Pearl of Great Price} % *PGP
	% % -----
		% \label{LAW-PGP}
	
		% \begin{tabular}{ll}
			% Total occurrences in the PGP & 
		% \end{tabular}
		
		% \begin{compactdesc}
			% \item[]
		% \end{compactdesc}
		
% ----------------------
\subsection{Key Phrases} % *KEYPHRASES *PHRASES
% ----------------------
	\label{LAW-PHRASES}

	\begin{compactdesc}
	
		\item[law of Christ] \textbf{7} | \hyperref[ROMANS-8-2]{Romans 8:2}; \hyperref[1CORINTHIANS-9-21]{1 Corinthians 9:21}; \hyperref[GALATIANS-6-2]{Galatians 6:2}; \hyperref[2NEPHI-25-27]{2 Nephi 25:27}; \hyperref[2NEPHI-26-1]{26:1}; \hyperref[DC88-21]{DC 88:21}; \hyperref[DC132-24]{DC 132:24} % law* christ*
			\begin{compactdesc}
				\item[Bible anchor verses] \phantom{.}
					\begin{compactdesc}
						\item[Galatians 6:2] \gr{ἀλλήλων τὰ βάρη βαστάζετε, καὶ οὕτως ἀναπληρώσετε \textbf{τὸν νόμον τοῦ Χριστοῦ}.}
					\end{compactdesc}
				\item[Commentary] See \hyperref[LAW-christ]{below}.
			\end{compactdesc}
			
		\item[law of death]
			
		\item[law of faith]
			\begin{compactdesc}
				\item[Bible anchor verses] \phantom{.}
					\begin{compactdesc}
						\item[Romans 3:27] 
					\end{compactdesc}
				% \item[Commentary] ---
			\end{compactdesc}
			
		\item[law of the gospel] See \hyperref[GOSPEL-PHRASES]{here}.
		
		\item[law of heaven] \textbf{1} | \hyperref[DC102-4]{DC 102:4}
			\begin{compactdesc}
				\item[Bible anchor verses] ---
				\item[Commentary] See also \hyperref[DC130-20]{DC 130:20}.
			\end{compactdesc}
			
		\item[law of sin]
		
		\item[spiritual law] \textbf{4} | \hyperref[ROMANS-7-14]{Romans 7:14}; \hyperref[ROMANS-8-2]{8:2}; \hyperref[2NEPHI-2-5]{2 Nephi 2:5}; \hyperref[DC29-34]{DC 29:34} % spirit* law*
			\begin{compactdesc}
				\item[Bible anchor verses] \phantom{.}
					\begin{compactdesc}
						\item[Romans 7:14] \gr{Οἴδαμεν γὰρ ὅτι \textbf{ὁ νόμος}} \hyperref[PNEUMATIKOS]{\grl{\textbf{πνευματικός}}} \gr{\textbf{ἐστιν}· ἐγὼ δὲ σάρκινός εἰμι, πεπραμένος ὑπὸ τὴν ἁμαρτίαν.}
						\item[Romans 8:2] \gr{\textbf{ὁ γὰρ νόμος τοῦ πνεύματος} τῆς ζωῆς ἐν Χριστῷ Ἰησοῦ ἠλευθέρωσέν σε ἀπὸ τοῦ νόμου τῆς ἁμαρτίας καὶ τοῦ θανάτου.}
					\end{compactdesc}
				% \item[Commentary] ---
			\end{compactdesc}
			
		\item[temporal law] \textbf{2} | \hyperref[2NEPHI-2-5]{2 Nephi 2:5}; \hyperref[DC29-34]{DC 29:34} % tempor* law*
		
	\end{compactdesc}
	
% % ----------------------
% \subsection{Bible Dictionaries} % *DICTIONARIES
% % ----------------------
	% \label{LAW-BD}

% ----------------------
\subsection{Restoration Usage} % *RESTORATION
% ----------------------
	\label{LAW-RESTORATION}

	% % -----
	% \subsubsection{Joseph Smith} % *JS
	% % -----
		% \label{LAW-JS}
	
		% \begin{compactdesc}
			% \item[{\hyperref[KEY]{Date/Source}}]
		% \end{compactdesc}
	
	% -----
	\subsubsection{Brigham Young} % *BY
	% -----
		\label{LAW-BY}
	
		\begin{compactdesc} % Format is: 21 May 1843 HC-a, or whatever date, whatever record-keeper (if there are multiple), then the passage
			\item[{\hyperref[EMS-1865-BY-8-OCT-1865]{8 October 1865}}] and now I don't wish to be at the trouble of quoting scripture passage after passage but suffice it to say \hl{when the celestial law is upon the earth the people believe that law and abide it they will be prepared and that law prepares them to inherit celestial glory in the presence of the Father and the Son who else will inherit it none of all the posterity of Adam upon the earth ever has been or ever will be none except those to whom the law of celestial kingdom is delivered and they that abide it can saved there in they will be crowned with glory immortality and eternal lives} now we differ a little form the Christian world with regard to this what will we say of them are they going to be thrust into hell yes will they dwell there forever and forever they will dwell somewhere unless they sin against the Holy Ghost will they become angels to the devil no they will not except they sin against the Holy Ghost will they reap the reward of those that do sin against the Holy Ghost no what will become of the inhabitants of the earth that have and now live all they that don't sin against the Holy Ghost will be saved in a kingdom somewhere and they fare better than ever imagined and Father Wesley get a greater glory [than] ever entered into his heart but [when/one/won't?] go where the Father and Son are and so it will be with others \hl{if the celestial law been in possession of John Wesley and reformers and they had abided the law of the celestial kingdom they would have become heirs to it otherwise they [can't?]}
		\end{compactdesc}
	
% ----------------------
\subsection{Commentary and Studies} % *COMMENTARY *STUDIES
% ----------------------
	\label{LAW-COMMENTARY}

	% % -----
	% \subsubsection{Key Points}
	% % -----
		% \label{LAW-KEYS}
	
		% \begin{compactitem}
			% \item
		% \end{compactitem}
		
	% -----
	\subsubsection{The ``law of Christ''}
	% -----
		\label{LAW-christ}
		
		See also BY, \href{https://drive.google.com/file/d/1wV-mRM932G6giTCQENRd8CTtRAMDGdxa/view?usp=sharing}{6 October 1867}: ``let that pure holy divine influence that is born in the heavens begotten of the Father given to the tabernacle placed within it let it rule prominent control until this sinful man is perfectly subdued and brought under subjection to the law of Christ''; and the \hyperref[TEMPLE-ENDOWMENT-ENDOW-PRE-1990-TELESTIAL]{endowment ceremony}, where Lucifer says, ``By what authority?'' and Peter responds, ``In the name of Jesus Christ, our Master.'' There is some relationship between the priesthood, Christ, and the law of Christ.
		
		See also BY, \href{https://drive.google.com/file/d/1LeTIp6sAiAgnwayXKIlXGZIXZ8CPetpO/view?usp=sharing}{19 June 1859}: ``if people would hearken to the whispering of that spirit they would be led in the paths of truth and righteousness if they would overcome the temptations of evil that is sin not in the spirit but in the flesh and let the spirit overcome the flesh this would bring into subjections the whole man to the law of Christ and make the children of man saints...we wish to see the kingdom of God and its fullness on the earth whose eyes do behold it will see a kingdom of purity a kingdom of holiness a people that is filled with the power of the upper world the power of God and sin will be overcome and this independent organization to mankind will be brought into subjection to that law we call it the law of Christ in [word/heart?] day [by] day it is the law of eternal life when we speak of law of Christ it is the pure and holy law that is calculated to keep matter in its organization go on the other hand you read of the first and second death what will be the result we witness it day by day the dissolution of the body but there is a second death still the other law [of the] gospel this law that is called the law of Christ the law of eternal lives that is set forth in Bible and Book of Mormon and all the revelations God gave to days of Adam until now let any man or people [follow?] out these laws it pertains to life it will not save their bodies from death that is the decree of Almighty but they will be made pure and holy and brought in celestial kingdom by this body being made pure by falling back to dust.'' 
		
		The \hyperref[LAW-BY]{BY quote above}, from 8 October 1865.
		
		See:
			\begin{compactitem}
				\item \hyperref[JOHN-5-26]{John 5:26}
				\item \hyperref[ROMANS-8-2]{Romans 8:2}
				\item \hyperref[GALATIANS-6-2]{Galatians 6:2}
				\item \hyperref[EPHESIANS-4-6]{Ephesians 4:6--16}
					\begin{compactitem}
						\item And the stuff after this is about the new birth, so there's clearly a relationship here. We're supposed to grow up to becoming like Christ.
					\end{compactitem}
				\item \hyperref[2NEPHI-9-5]{2 Nephi 9:5}---is it part of the law of Christ that he has to allow himself to be subject to man, so that others will be subject to him? See also \hyperref[3NEPHI-27-14]{3 Nephi 27:14} for a related idea.
				\item \hyperref[2NEPHI-26-1]{2 Nephi 26:1}
					\begin{compactitem}
						\item See also \hyperref[3NEPHI-18-14]{3 Nephi 18:14} and probably other related verses.
					\end{compactitem}
				\item \hyperref[2NEPHI-32-6]{2 Nephi 32:6} (and \hyperref[2NEPHI-26-1]{2 Nephi 26:1})
				\item \hyperref[3NEPHI-11]{3 Nephi 11}
					\begin{compactitem}
						\item Starting around 31. I would say, see if you can make a connection between the law of Christ and the ``doctrine'' of Christ, as explained here and at 2 Nephi 31 and in other places. Maybe also connect to definitions of the gospel.
						\item Also 3 Nephi 12:19--20?
					\end{compactitem}
				\item \hyperref[3NEPHI-15-9]{3 Nephi 15:9--10}
				\item \hyperref[DC38-22]{DC 38:22}
				\item \hyperref[DC38-32]{DC 38:32}
				\item \hyperref[DC45-59]{DC 45:59}
				\item \hyperref[DC51-15]{DC 51:15} (also mentioned earlier in the section)
				\item \hyperref[DC88-5]{DC 88:5--13}
				\item \hyperref[DC105-4]{DC 105:4--5}
				\item \hyperref[DC107-3]{DC 107:3}
				\item \hyperref[DC132-24]{DC 132:24}
			\end{compactitem}